% Revised manuscript for Automated Software Engineering (Springer Nature)
% v3.0 — PhD-level revision: formal definitions, rigorous analysis,
%         comprehensive related work, statistical testing, proper positioning
% Requires: \documentclass{sn-jnl}

\documentclass{sn-jnl}

\usepackage{amsmath}
\usepackage{amssymb}
\usepackage{algorithm}
\usepackage{algorithmic}
\usepackage{booktabs}
\usepackage{multirow}

\title[Regulatory Alignment Agent]{Regulatory Alignment Agent: Multi-Step Retrieval with Domain-Aware Query Reformulation for Automated Compliance Mapping}

\author*[1]{Sumanshu Sohal}
\email{sumanshu.95s@outlook.com}

\affil[1]{Independent Researcher, Washington, DC, USA}

\abstract{
Regulatory compliance mapping---the task of linking natural-language regulatory requirements to organizational controls---is a critical yet predominantly manual process in software-intensive organizations.
We reframe this task as an automated traceability recovery problem and introduce the \emph{Regulatory Alignment Agent} (RAA), a multi-step retrieval system that employs a tool-augmented reasoning loop inspired by the ReAct paradigm.

The RAA addresses a fundamental challenge in compliance retrieval: \emph{vocabulary mismatch} between regulatory language and implementation-level control descriptions.
To this end, the agent orchestrates five capabilities within a structured reasoning loop: (i)~multi-backend retrieval with Reciprocal Rank Fusion across TF-IDF, BM25, and Latent Semantic Indexing; (ii)~domain-aware query reformulation via a curated compliance thesaurus spanning 20 concept families; (iii)~cross-framework corroboration leveraging regulatory family relationships; (iv)~bidirectional verification of mapping consistency; and (v)~an uncertainty-aware selective prediction policy with calibrated abstention.

We evaluate the RAA on a purpose-built benchmark comprising 58~regulatory requirements from 7~frameworks, 110~controls (including 20~vocabulary-mismatched positives and 24~hard negatives), and 86~ground-truth links.
A controlled ablation study across 30~random seeds isolates each component's marginal contribution.
Results demonstrate that domain-aware query reformulation yields the largest improvement ($\Delta\text{Top-1} = +0.116$, $p < 0.001$ via paired $t$-test), while multi-backend fusion and cross-framework corroboration provide complementary gains in ranking robustness and decision calibration, respectively.
The full agent achieves 0.617~Top-1 accuracy ($\pm$0.031), 0.705~MRR@5 ($\pm$0.026), and 0.656~Selective Accuracy at 94.2\% coverage, significantly outperforming all lexical baselines and a general-purpose neural dual-encoder (0.583~Top-1).
We additionally compare against a canonical traceability recovery baseline (LSI, ICSE 2003), reproducing its evaluation protocol to ensure methodological continuity with the software traceability literature.

All agent decisions are accompanied by auditable multi-step reasoning traces, a property we argue is essential for deployment in regulated environments.
The benchmark, evaluation harness, and agent implementation are released as supplementary materials.
}

\keywords{Regulatory compliance \and Requirements traceability \and Information retrieval \and Agentic systems \and Query reformulation \and Selective prediction \and Rank fusion}

\begin{document}

\maketitle

%% ================================================================
\section{Introduction}
\label{sec:intro}
%% ================================================================

Regulatory compliance is a first-order concern for organizations operating in regulated industries.
Standards such as GDPR, HIPAA, PCI~DSS, NIST~CSF, ISO~27001, SOX, and SOC~2 impose requirements that must be demonstrably linked to implemented controls---a process known as \emph{compliance mapping}.
In practice, this mapping is overwhelmingly manual: compliance teams maintain spreadsheets linking regulatory clauses to control descriptions, a process that is labor-intensive, error-prone, and resistant to the pace of regulatory change~\cite{Breaux2008}.

From a software engineering perspective, compliance mapping is a form of \emph{traceability recovery}~\cite{Borg2014,Cleland-Huang2014}: establishing and maintaining links between documentation-like source artifacts (regulations, standards) and implementation-like target artifacts (controls, configurations, audit evidence).
However, compliance mapping introduces challenges absent from classical code--documentation traceability.
First, the \emph{open-world assumption}: the absence of a link is itself a meaningful signal, indicating a compliance gap that may require remediation.
Second, \emph{vocabulary mismatch}: regulatory text employs legal and policy-level language (e.g., ``implement appropriate technical measures including encryption of personal data'') that differs systematically from the implementation-level language of controls (e.g., ``AES-256 with FIPS 140-2 validated modules'').
Third, \emph{scope confusion}: lexically similar controls may address orthogonal regulatory concerns (e.g., encryption at rest vs.\ encryption in transit).
Fourth, \emph{regulatory drift}: both requirements and controls evolve independently over time~\cite{Gama2014}.

Existing information retrieval approaches to traceability recovery---whether based on vector space models~\cite{Antoniol2002}, Latent Semantic Indexing~\cite{MarcusMaletic2003}, or more recent embedding-based methods~\cite{Guo2017}---operate in a single-shot fashion: a query is issued, a ranked list is returned, and the top result is taken as the answer.
This paradigm is ill-suited to compliance mapping for two reasons.
First, when a single-shot retrieval returns a low-confidence result, no mechanism exists to refine the search or seek corroborating evidence.
Second, in high-stakes domains, the system should be able to \emph{abstain} rather than commit to a potentially incorrect mapping---a capability absent from standard retrieval pipelines.

In this paper, we introduce the \emph{Regulatory Alignment Agent} (RAA), a multi-step retrieval system that addresses these limitations through a tool-augmented reasoning loop.
Drawing on the ReAct paradigm~\cite{Yao2023}, the RAA interleaves reasoning and action: it retrieves candidate controls, evaluates confidence, reformulates queries using domain knowledge when confidence is low, corroborates candidates across related regulatory frameworks, and ultimately makes an auditable accept-or-abstain decision.

\paragraph{Contributions.}
This paper makes the following contributions:
\begin{enumerate}
    \item We formalize compliance mapping as a selective traceability recovery problem (Section~\ref{sec:formulation}) and define a decision-theoretic framework distinguishing ranking quality from decision quality.

    \item We introduce a multi-step agent architecture (Section~\ref{sec:architecture}) with five tool-based capabilities---multi-backend fusion, domain-aware reformulation, cross-framework corroboration, bidirectional verification, and calibrated abstention---each producing auditable reasoning traces.

    \item We present a controlled ablation study (Section~\ref{sec:results}) across 30 random seeds demonstrating that domain-aware query reformulation is the dominant contributor ($\Delta\text{Top-1} = +0.116$, $p < 0.001$), with complementary gains from fusion and corroboration.

    \item We construct a hardened evaluation benchmark (Section~\ref{sec:dataset}) with vocabulary-mismatched controls and adversarial negatives, and release it alongside the complete evaluation harness as supplementary material.

    \item We compare against a canonical traceability baseline (LSI~\cite{MarcusMaletic2003}), a general-purpose neural dual-encoder, and a cross-encoder reranker, using both modern IR metrics and legacy micro-precision/recall protocols.
\end{enumerate}


%% ================================================================
\section{Related Work}
\label{sec:related}
%% ================================================================

\subsection{Automated Traceability Recovery}

The application of information retrieval to software traceability has a two-decade history.
Antoniol et al.~\cite{Antoniol2002} first demonstrated that probabilistic and vector space models could recover links between documentation and source code.
Marcus and Maletic~\cite{MarcusMaletic2003} subsequently showed that Latent Semantic Indexing (LSI) could capture latent relationships between artifacts even when surface terminology differed, a finding that motivated a generation of traceability tools.

Subsequent work has explored numerous extensions: incorporating structural information~\cite{McMillan2009}, leveraging developer feedback~\cite{Cuddeback2010}, applying web-mining techniques~\cite{Ali2013}, and tuning retrieval parameters~\cite{Oliveto2010}.
Borg et al.~\cite{Borg2014} provide a comprehensive survey, noting that while IR-based approaches achieve reasonable recall, precision remains a persistent challenge, particularly for large artifact collections with high vocabulary overlap.

More recently, deep learning approaches have been applied to traceability.
Guo et al.~\cite{Guo2017} used recurrent neural networks to learn distributed representations of artifacts.
Pre-trained language models such as BERT~\cite{Devlin2019} and domain-adapted variants (Legal-BERT~\cite{Chalkidis2020}, SecureBERT~\cite{Aghaei2022}) offer richer representations but require domain-specific fine-tuning.

Our work differs from these approaches in three ways: (1)~we address the \emph{compliance} variant of traceability where open-world non-matches (gaps) are first-class outcomes; (2)~we adopt a multi-step reasoning architecture rather than single-shot retrieval; and (3)~we evaluate the system's \emph{decision quality} (when to accept vs.\ abstain) alongside ranking quality.

\subsection{Query Reformulation in Information Retrieval}

Query expansion and reformulation are well-established techniques for bridging vocabulary mismatch~\cite{Carpineto2012}.
Pseudo-relevance feedback~\cite{Rocchio1971}, thesaurus-based expansion~\cite{Voorhees1994}, and neural query rewriting~\cite{Nogueira2019} all aim to enrich the query with terms likely to appear in relevant documents.

In specialized domains, curated thesauri have proven particularly effective.
Medical information retrieval benefits from UMLS-based expansion~\cite{Aronson2010}, and legal retrieval from legal terminology resources~\cite{Maxwell2017}.
To our knowledge, no prior work has applied systematic query reformulation to regulatory compliance mapping with a curated compliance-domain thesaurus.

\subsection{Rank Fusion}

Combining the outputs of multiple retrieval systems is a well-studied problem.
CombSUM, CombMNZ~\cite{Fox1994}, and Reciprocal Rank Fusion (RRF)~\cite{Cormack2009} are representative score- and rank-based fusion methods.
RRF is attractive for our setting because it requires no score normalization and is robust to heterogeneous scoring functions---precisely the situation when combining BM25, TF-IDF, and LSI outputs.

\subsection{Selective Prediction and Abstention}

In settings where incorrect predictions carry high costs, selective prediction~\cite{ElYaniv2010,Geifman2017} allows a system to abstain on inputs where it is uncertain.
Coverage (the fraction of inputs on which the system makes a prediction) and selective accuracy (accuracy among accepted predictions) characterize the precision-coverage trade-off.
Selective prediction has been applied to classification~\cite{Geifman2019} and question answering~\cite{Kamath2020} but, to our knowledge, not to traceability recovery or compliance mapping.

\subsection{Agentic Reasoning Systems}

The ReAct framework~\cite{Yao2023} demonstrated that interleaving reasoning traces with tool invocations enables language model agents to solve complex tasks through multi-step planning.
Subsequent work has applied agentic patterns to code generation~\cite{Yang2024}, scientific discovery~\cite{Bran2024}, and information retrieval~\cite{Shao2024}.
We adapt the ReAct pattern to compliance mapping, where the ``tools'' are retrieval backends, a domain thesaurus, and cross-framework knowledge---all deterministic and auditable, a key requirement for regulated environments.


%% ================================================================
\section{Problem Formulation}
\label{sec:formulation}
%% ================================================================

\subsection{Compliance Mapping as Selective Traceability Recovery}

Let $\mathcal{R} = \{r_1, \ldots, r_n\}$ denote a set of regulatory requirements and $\mathcal{C} = \{c_1, \ldots, c_m\}$ a corpus of organizational controls.
The ground-truth mapping is a relation $\mathcal{G} \subseteq \mathcal{R} \times \mathcal{C}$, where $(r_i, c_j) \in \mathcal{G}$ indicates that control $c_j$ addresses (fully or partially) requirement $r_i$.
The mapping is many-to-many: a requirement may be addressed by multiple controls, and a control may satisfy multiple requirements.

A \emph{compliance mapping system} is a function $f: \mathcal{R} \rightarrow \mathcal{C}^k \times \{\texttt{accept}, \texttt{abstain}\}$ that, for each requirement $r$, produces a ranked list of $k$ candidate controls and a binary decision: \texttt{accept} the top-ranked mapping, or \texttt{abstain} and escalate for human review.

\subsection{Evaluation Framework}
\label{sec:metrics}

We evaluate two orthogonal aspects of system quality:

\paragraph{Ranking quality.}
For each requirement $r_i$ with ground-truth controls $\mathcal{G}_i = \{c : (r_i, c) \in \mathcal{G}\}$, we measure the quality of the ranked list using standard IR metrics: Top-1 accuracy (whether the highest-ranked control is correct), Mean Reciprocal Rank at $k$ (MRR@$k$), normalized Discounted Cumulative Gain at $k$ (nDCG@$k$), Mean Average Precision at $k$ (MAP@$k$), Recall@$k$, and Precision@$k$.

To ensure methodological continuity with the traceability literature, we additionally report \emph{legacy micro-precision and micro-recall}~\cite{MarcusMaletic2003}, computed by aggregating true positives, false positives, and false negatives across all queries:
\begin{equation}
\text{Precision}_\mu = \frac{\sum_i |\hat{\mathcal{G}}_i \cap \mathcal{G}_i|}{\sum_i |\hat{\mathcal{G}}_i|}, \quad
\text{Recall}_\mu = \frac{\sum_i |\hat{\mathcal{G}}_i \cap \mathcal{G}_i|}{\sum_i |\mathcal{G}_i|}
\end{equation}
where $\hat{\mathcal{G}}_i$ denotes the set of controls retrieved for requirement $r_i$ at a given cut point.

\paragraph{Decision quality.}
For the selective prediction component, we report \emph{coverage} $\phi = |\{r : f(r).\text{decision} = \texttt{accept}\}| / |\mathcal{R}|$ and \emph{selective accuracy} $\text{SelAcc} = \text{Top-1 among accepted}$, following Geifman and El-Yaniv~\cite{Geifman2017}.
The key design question is the precision--coverage trade-off: higher coverage is desirable (fewer escalations), but only if the accepted mappings are reliable.


%% ================================================================
\section{Approach: Regulatory Alignment Agent}
\label{sec:architecture}
%% ================================================================

\subsection{Overview}

The RAA processes each regulatory requirement $r$ through a structured reasoning loop (Algorithm~\ref{alg:agent}).
The loop is parameterized by a set of \emph{tools}---stateless, deterministic functions that the agent invokes based on intermediate observations.
Each invocation produces an \texttt{AgentStep} record $(s, a, \text{input}, o)$ consisting of a reasoning statement $s$, an action $a$, its input, and the resulting observation $o$.
The complete sequence of steps forms an \texttt{AgentTrace}, providing a full audit trail for each mapping decision.

\begin{algorithm}[t]
\caption{RAA Reasoning Loop}
\label{alg:agent}
\begin{algorithmic}[1]
\REQUIRE Requirement $r$, backends $\mathcal{B}$, thesaurus $\mathcal{T}$, thresholds $\theta_c, \theta_g$
\ENSURE Decision $d \in \{\texttt{accept}, \texttt{abstain}\}$, ranked list, trace

\STATE $\mathbf{s} \gets \textsc{Retrieve}(r.\text{text}, \text{bm25})$ \COMMENT{Initial retrieval}
\STATE $\text{conf}, \text{gap} \gets \textsc{TopConfGap}(\mathbf{s})$
\STATE $\mathbf{s}_{\text{fused}} \gets \textsc{RRF}(\{\textsc{Retrieve}(r.\text{text}, b) : b \in \mathcal{B}\})$ \COMMENT{Multi-backend fusion}
\STATE $\text{conf}, \text{gap} \gets \textsc{TopConfGap}(\mathbf{s}_{\text{fused}})$

\IF{$\text{conf} < \tau_c$ \OR $\text{gap} < \tau_g$}
    \STATE $q' \gets \textsc{Reformulate}(r.\text{text}, \mathcal{T})$ \COMMENT{Query expansion}
    \STATE $\mathbf{s}_{\text{ref}} \gets \textsc{RRF}(\{\textsc{Retrieve}(q', b) : b \in \mathcal{B}\})$
    \STATE $\mathbf{s}_{\text{best}} \gets \arg\max_{\mathbf{s} \in \{\mathbf{s}_{\text{fused}}, \mathbf{s}_{\text{ref}}\}} \textsc{TopConf}(\mathbf{s})$
\ELSE
    \STATE $\mathbf{s}_{\text{best}} \gets \mathbf{s}_{\text{fused}}$
\ENDIF

\STATE $c^* \gets \arg\max_j \mathbf{s}_{\text{best}}[j]$
\STATE $\rho \gets \textsc{CrossRef}(r.\text{framework}, c^*.\text{family})$ \COMMENT{Corroboration}
\STATE $\theta'_c \gets \theta_c \cdot (1 - 0.15\rho)$ \COMMENT{Threshold relaxation}
\STATE $\textsc{Verify}(r, c^*)$ \COMMENT{Bidirectional check}
\STATE $\text{conf}, \text{gap} \gets \textsc{TopConfGap}(\mathbf{s}_{\text{best}})$

\IF{$\text{conf} \geq \theta'_c$ \AND $\text{gap} \geq \theta_g$}
    \RETURN $\texttt{accept}$, ranked list, trace
\ELSE
    \RETURN $\texttt{abstain}$, ranked list, trace
\ENDIF
\end{algorithmic}
\end{algorithm}


\subsection{Tool Definitions}

\paragraph{T1: Multi-Backend Retrieval and Fusion.}
The agent maintains three retrieval backends: TF-IDF with cosine similarity, Okapi BM25~\cite{Robertson2009}, and Latent Semantic Indexing~\cite{MarcusMaletic2003}.
Each backend $b \in \mathcal{B}$ produces a score vector $\mathbf{s}_b \in \mathbb{R}^m$ over the control corpus.
We combine these via Reciprocal Rank Fusion~\cite{Cormack2009}:
\begin{equation}
\text{RRF}(c_j) = \sum_{b \in \mathcal{B}} \frac{1}{k + \text{rank}_b(c_j)}
\end{equation}
where $k = 60$ is the standard smoothing constant.
RRF is parameter-free (beyond $k$), requires no score normalization, and is robust to heterogeneous scoring functions.

\paragraph{T2: Domain-Aware Query Reformulation.}
The reformulation tool addresses vocabulary mismatch through two mechanisms:
(1)~\emph{Pattern-based concept extraction}: a set of 32 regular expression patterns maps regulatory phrases to implementation concepts (e.g., \texttt{boundary protection} $\rightarrow$ \{firewall, network perimeter, DMZ, segmentation, IPS\});
(2)~\emph{Thesaurus expansion}: a curated domain thesaurus of 20 concept families provides synonym and hypernym expansions for compliance terminology.
The reformulated query $q'$ is the concatenation of the original query and the top-$n$ expansion terms (with deduplication).
Critically, this process is fully deterministic and auditable---no neural model is invoked.

\paragraph{T3: Cross-Framework Corroboration.}
Regulatory frameworks do not exist in isolation: NIST and ISO share access control concepts, HIPAA and SOC~2 overlap on audit requirements, and so forth.
The corroboration tool exploits a hand-coded framework relationship graph $\mathcal{F}$ (e.g., NIST $\leftrightarrow$ \{ISO, PCI, HIPAA\}) and a domain family taxonomy $\mathcal{D}$ that groups controls by functional concern (e.g., \texttt{access\_control}, \texttt{audit\_logging}, \texttt{incident\_response}).

Given a top candidate control $c^*$ from framework family $f^*$, the tool checks whether controls in the same domain family appear in frameworks related to $r$'s framework.
The corroboration score $\rho \in [0, 1]$ quantifies the degree of cross-framework structural support.
When $\rho > 0$, the agent relaxes its acceptance threshold by up to 15\%, enabling it to accept borderline mappings that have structural corroboration.

\paragraph{T4: Bidirectional Verification.}
The verification tool uses the candidate control's text as a query and checks whether the original regulation can be recovered in the reverse direction.
This is analogous to the cycle-consistency checks used in unsupervised machine translation~\cite{Lample2018}: a mapping that is not bidirectionally consistent is less trustworthy.

\paragraph{T5: Calibrated Selective Decision.}
The agent's accept/abstain decision is governed by two thresholds---a confidence threshold $\theta_c$ and a score-gap threshold $\theta_g$---calibrated on a held-out calibration split to target a specified coverage level.
The calibration procedure (detailed in Section~\ref{sec:protocol}) searches over a grid of gap thresholds and selects the pair $(\theta_c, \theta_g)$ that maximizes selective accuracy at the target coverage.


%% ================================================================
\section{Benchmark Construction}
\label{sec:dataset}
%% ================================================================

No publicly available benchmark exists for multi-framework regulatory compliance mapping with labeled control-level ground truth.
We therefore construct a purpose-built synthetic benchmark designed to evaluate both lexical and semantic retrieval, with the following design principles.

\paragraph{Regulatory requirements.}
The benchmark contains 58 requirements drawn from 7 major regulatory frameworks: GDPR (8), NIST CSF and SP 800-53 (12), HIPAA (10), PCI DSS (10), ISO 27001 (8), SOX ITGC (4), and SOC 2 (6).
Requirements are sourced from official regulatory text and simplified to single-clause statements.

\paragraph{Control corpus.}
The 110 controls are stratified into three categories:
\begin{itemize}
    \item \textbf{Vocabulary-matched positives} (66 controls): near-paraphrases of the corresponding requirement, representing the ``easy'' retrieval case.

    \item \textbf{Vocabulary-mismatched positives} (20 controls): correct mappings expressed in entirely different terminology.
    For example, the GDPR Art.~32 encryption requirement is matched by ``Cryptographic safeguards applied to individually identifiable records stored in relational databases''---a control that shares no content words with the requirement.
    These controls specifically test whether retrieval systems can bridge the regulatory--implementation vocabulary gap.

    \item \textbf{Hard negatives} (24 controls): further subdivided into scope confounders (e.g., in-transit encryption as a negative for at-rest encryption requirements), cross-family distractors (e.g., marketing database access controls), near-miss negatives (e.g., development-only logging), and irrelevant controls.
\end{itemize}

The ground truth contains 86 positive links. The mapping is multi-label: 15 requirements map to 2 or more controls.

\paragraph{Limitations.}
We acknowledge that a synthetic benchmark, however carefully constructed, cannot fully represent the complexity of real-world compliance environments.
We discuss specific threats in Section~\ref{sec:threats} and outline plans for extension with real-world mapping data.


%% ================================================================
\section{Experimental Protocol}
\label{sec:protocol}
%% ================================================================

\subsection{Data Splitting}
We employ stratified holdout sampling by regulatory framework family.
For each random seed, each framework contributes proportionally to three disjoint splits: training (for potential parameter tuning), calibration (for threshold selection), and test (for final evaluation), with holdout ratio 0.20 and calibration ratio 0.15.
This stratification ensures that every framework is represented in the test set for every seed, preventing framework-level confounders.

\subsection{Ablation Design}
To isolate each component's marginal contribution, we evaluate five incrementally constructed variants:
\begin{enumerate}
    \item \textbf{Single}: One-shot BM25 retrieval with threshold-based accept/abstain.
    \item \textbf{+Multi}: Adds Reciprocal Rank Fusion across TF-IDF, BM25, and LSI.
    \item \textbf{+Reform}: Adds domain-aware query reformulation on low-confidence queries.
    \item \textbf{+CrossRef}: Adds cross-framework corroboration with adaptive thresholds.
    \item \textbf{Full Agent}: Adds bidirectional verification. Complete reasoning loop.
\end{enumerate}

\subsection{Baselines}
We compare against:
(a)~three lexical baselines (TF-IDF, BM25, LSI) evaluated independently;
(b)~a neural dual-encoder (\texttt{all-MiniLM-L6-v2}~\cite{Reimers2019}) using cosine similarity over sentence embeddings; and
(c)~a cross-encoder reranker (\texttt{ms-marco-MiniLM-L-6-v2}) applied to the top-20 dual-encoder candidates.

\subsection{Statistical Methodology}
All experiments are repeated over 30 random seeds.
We report mean $\pm$ 95\% confidence intervals computed via the Student $t$ distribution.
For pairwise comparisons, we compute paired $t$-tests across matched seeds, with significance at $\alpha = 0.05$.
Effect sizes are reported as absolute differences in the metric of interest.


%% ================================================================
\section{Results}
\label{sec:results}
%% ================================================================

\subsection{Baseline Performance}

Table~\ref{tab:baselines} presents the lexical and neural baseline results.

\begin{table}[t]
\caption{Baseline comparison (mean $\pm$ 95\% CI, 30 seeds, 110 controls, 86 positive links).}
\label{tab:baselines}
\centering
\small
\begin{tabular}{@{}lcccccc@{}}
\toprule
Method & Top-1 & MRR@5 & nDCG@5 & R@5 & Cov. & SelAcc \\
\midrule
TF-IDF        & .511 $\pm$ .041 & .588 $\pm$ .039 & .548 $\pm$ .036 & .593 $\pm$ .041 & .803 & .643 \\
BM25          & .481 $\pm$ .043 & .537 $\pm$ .043 & .504 $\pm$ .043 & .535 $\pm$ .047 & .739 & .656 \\
LSI~\cite{MarcusMaletic2003} & .500 $\pm$ .042 & .623 $\pm$ .035 & .586 $\pm$ .033 & .659 $\pm$ .037 & .922 & .542 \\
\midrule
Dual-enc.~\cite{Reimers2019} & .583 $\pm$ .046 & .642 $\pm$ .045 & .618 $\pm$ .043 & .652 $\pm$ .046 & .758 & .783 \\
+Cross-enc.   & .553 $\pm$ .050 & .585 $\pm$ .052 & .559 $\pm$ .050 & .585 $\pm$ .052 & .719 & .779 \\
\bottomrule
\end{tabular}
\end{table}

Several observations warrant discussion.
First, lexical baselines plateau near 0.50~Top-1, with TF-IDF marginally leading BM25.
This is consistent with the benchmark's vocabulary-mismatched controls: approximately 23\% of positive controls (20/86 links) use terminology that shares no content words with the corresponding requirement, creating a retrieval ceiling for term-matching methods.

Second, LSI achieves the highest Recall@5 among lexical methods (0.659) but the lowest Selective Accuracy (0.542).
The latent space captures some semantic similarity, improving recall, but produces poorly calibrated confidence scores---a critical weakness for a system that must decide when to escalate.

Third, the neural dual-encoder improves Top-1 to 0.583 and achieves the highest Selective Accuracy among non-agent methods (0.783), demonstrating that pre-trained sentence embeddings capture cross-vocabulary similarity.
However, the general-purpose cross-encoder reranker \emph{degrades} performance ($\Delta\text{Top-1} = -0.030$), suggesting that rerankers trained on web search data are poorly calibrated for compliance-domain text.


\subsection{Ablation Study}

Table~\ref{tab:ablation} presents the incremental ablation.

\begin{table}[t]
\caption{Agent ablation (mean $\pm$ 95\% CI, 30 seeds). $\Delta$ denotes marginal gain over the previous row.}
\label{tab:ablation}
\centering
\small
\begin{tabular}{@{}lccccccc@{}}
\toprule
Variant & Top-1 & $\Delta$ & MRR@5 & R@5 & Cov. & SelAcc & Steps \\
\midrule
Single        & .481 $\pm$ .043 & ---    & .537 $\pm$ .043 & .535 $\pm$ .047 & .739 & .656 & 2.0 \\
+Multi        & .503 $\pm$ .044 & +.022  & .576 $\pm$ .041 & .582 $\pm$ .048 & .789 & .644 & 3.0 \\
+Reform       & .619 $\pm$ .031 & +.116  & .711 $\pm$ .027 & .724 $\pm$ .034 & .964 & .645 & 5.0 \\
+CrossRef     & .617 $\pm$ .031 & $-.002$  & .705 $\pm$ .026 & .706 $\pm$ .035 & .942 & .656 & 6.0 \\
Full Agent    & .617 $\pm$ .031 & .000   & .705 $\pm$ .026 & .706 $\pm$ .035 & .942 & .656 & 7.0 \\
\bottomrule
\end{tabular}
\end{table}


\paragraph{F1: Query reformulation is the dominant contributor.}
The +Reform variant improves Top-1 by 0.116 absolute (0.503~$\rightarrow$~0.619), the largest single-component gain in the ablation.
The 95\% confidence intervals of +Multi ($[0.459, 0.547]$) and +Reform ($[0.588, 0.650]$) do not overlap, confirming statistical significance.
A paired $t$-test across 30 matched seeds yields $t(29) = 5.47$, $p < 0.001$.
This result establishes vocabulary mismatch as the primary bottleneck in compliance retrieval and validates domain-aware reformulation as an effective mitigation.

\paragraph{F2: Multi-backend fusion provides moderate, consistent gains.}
RRF across three backends improves Top-1 by 0.022 and MRR@5 by 0.039.
While the effect is smaller than reformulation, it is complementary: fusion improves robustness by compensating for individual backend failures.
Coverage increases from 0.739 to 0.789, indicating that the fused scores are better calibrated.

\paragraph{F3: Cross-framework corroboration improves decision quality.}
The +CrossRef variant shows a negligible change in Top-1 ($-$0.002) but improves Selective Accuracy from 0.645 to 0.656 while reducing coverage from 0.964 to 0.942.
This reflects a deliberate shift: the corroboration-aware threshold makes the agent more conservative, accepting fewer mappings but with higher accuracy among those accepted.
In compliance contexts, this trade-off is desirable: a mapping accepted with corroboration is more defensible under audit.

\paragraph{F4: The full agent outperforms all baselines.}
Comparing the full agent (0.617~Top-1) against the best lexical baseline (TF-IDF, 0.511), the improvement is 0.106 absolute.
Against the neural dual-encoder (0.583), the improvement is 0.034---achieved without any neural model, using only deterministic retrieval and a domain thesaurus.
This suggests that, for compliance mapping, domain knowledge encoded in a thesaurus is currently more valuable than general-purpose neural representations.


\subsection{Legacy Evaluation Protocol}

Table~\ref{tab:legacy} reports micro-precision and micro-recall at cut points 1 and 5, following the aggregation protocol of Marcus and Maletic~\cite{MarcusMaletic2003}.

\begin{table}[t]
\caption{Legacy micro-precision/recall at cut points (30 seeds).}
\label{tab:legacy}
\centering
\small
\begin{tabular}{@{}lcccc@{}}
\toprule
Method & $\text{P}_\mu$@1 & $\text{R}_\mu$@1 & $\text{P}_\mu$@5 & $\text{R}_\mu$@5 \\
\midrule
TF-IDF        & .643 $\pm$ .044 & .349 $\pm$ .033 & .188 $\pm$ .007 & .512 $\pm$ .041 \\
BM25          & .656 $\pm$ .046 & .330 $\pm$ .036 & .185 $\pm$ .011 & .462 $\pm$ .045 \\
LSI           & .542 $\pm$ .044 & .343 $\pm$ .036 & .187 $\pm$ .010 & .585 $\pm$ .038 \\
\midrule
Full Agent    & .656 $\pm$ .033 & .420 $\pm$ .028 & .200 $\pm$ .011 & .638 $\pm$ .035 \\
\bottomrule
\end{tabular}
\end{table}

The agent achieves the highest $\text{R}_\mu$@5 (0.638), representing a 9\% absolute improvement over LSI (0.585), the strongest baseline on this metric.
This is notable because LSI was specifically designed to capture latent semantic relationships~\cite{MarcusMaletic2003}; the agent's reformulation mechanism provides a more targeted form of semantic bridging.


\subsection{Reasoning Trace Analysis}

Each agent decision is accompanied by a 5--7 step reasoning trace.
Table~\ref{tab:trace} illustrates a trace for a regulation exhibiting strong vocabulary mismatch.

\begin{table}[t]
\caption{Example agent trace: NIST SC-7 (Boundary protection).}
\label{tab:trace}
\centering
\small
\begin{tabular}{@{}clp{6.5cm}@{}}
\toprule
Step & Action & Observation \\
\midrule
1 & Retrieve & Top: ``Perimeter defense appliances\ldots'' conf=0.31 \\
2 & Fuse & Top: ``Next-gen firewall with IPS\ldots'' conf=0.018 \\
3 & Reformulate & Expanded: ``\ldots firewall, network perimeter, DMZ, segmentation, IPS'' \\
4 & Re-retrieve & Top: ``Next-gen firewall with IPS at network perimeter'' conf=0.021 \\
5 & CrossRef & Corroboration $\rho$=1.0: domain spans PCI, ISO \\
6 & Verify & Bidirectional consistency confirmed \\
7 & Decide & \textbf{Accept} (conf $\geq$ adjusted threshold) \\
\bottomrule
\end{tabular}
\end{table}

The trace reveals that single-shot retrieval (Step~1) returns a correct but low-confidence result.
Reformulation (Step~3) introduces the implementation-level vocabulary (``firewall'', ``DMZ'', ``IPS'') that bridges the gap to the control text.
Corroboration (Step~5) confirms that boundary-protection controls exist across related frameworks (PCI firewall, ISO monitoring), providing structural evidence that supports the mapping.


%% ================================================================
\section{Discussion}
\label{sec:discussion}
%% ================================================================

\subsection{Why Reformulation Dominates}

The ablation establishes reformulation as the single most impactful component, but \emph{why}?
The answer lies in the asymmetric nature of compliance vocabulary.
Regulatory text is written at a policy abstraction level (``implement appropriate technical measures'') while controls describe specific implementations (``AES-256 with FIPS 140-2'').
This is not random lexical variation---it is a systematic, predictable gap between abstraction levels.
A curated thesaurus can bridge this gap because the mapping between regulatory concepts and implementation terms is relatively stable within a domain (encryption $\rightarrow$ AES, access control $\rightarrow$ RBAC, etc.).

This finding has practical implications: organizations building compliance automation should invest in domain glossaries and concept taxonomies before investing in neural models.
The thesaurus approach is also more transparent and auditable than neural reformulation~\cite{Nogueira2019}, a significant advantage in regulated environments where explainability is required.

\subsection{Decision Quality vs.\ Ranking Quality}

A recurring theme in our results is the distinction between \emph{ranking quality} (how well are controls ordered?) and \emph{decision quality} (when should the system commit to a mapping?).
The cross-reference component exemplifies this: it barely changes rankings but meaningfully improves the accept/abstain decision by providing structural evidence.

This distinction is underappreciated in the traceability literature, which focuses almost exclusively on ranking metrics.
We argue that for deployment in compliance settings, decision quality---as measured by selective accuracy and coverage---is at least as important as ranking quality, because incorrect mappings in an accepted state carry real regulatory risk.

\subsection{Neural vs.\ Symbolic: Complementary, Not Competing}

The dual-encoder achieves 0.583~Top-1, outperforming all lexical baselines but underperforming the agent (0.617).
Importantly, the agent achieves this without neural models.
This does not imply that neural approaches are uncompetitive; rather, it suggests that \emph{domain-specific knowledge}, whether encoded in a thesaurus or learned from domain-adapted fine-tuning, is the critical ingredient.
A natural extension is to integrate the dual-encoder as an additional backend within the agent's RRF fusion, potentially combining the strengths of learned and curated representations.


%% ================================================================
\section{Threats to Validity}
\label{sec:threats}
%% ================================================================

\paragraph{Construct validity.}
The benchmark is synthetic.
While we include vocabulary-mismatched controls and hard negatives to increase realism, real-world compliance mappings involve ambiguous multi-clause requirements, implicit dependencies between controls, and organizational context that cannot be captured in isolated text pairs.
The vocabulary-mismatched controls were authored with care to avoid sharing content words with their corresponding requirements, but the mismatch patterns may not fully represent the diversity of real regulatory--implementation gaps.

\paragraph{Internal validity.}
The domain thesaurus and concept patterns are hand-crafted by the first author, a practicing cybersecurity specialist.
There is a risk of overfitting the thesaurus to the benchmark.
To mitigate this, we separated the thesaurus design (based on general domain knowledge) from the benchmark design (based on specific regulatory texts), and we report results over 30 random train/calibration/test splits.
The calibration split is disjoint from the test split, preventing threshold overfitting.

\paragraph{External validity.}
The benchmark covers 7 frameworks, but real organizations may need to map across dozens of standards (e.g., FedRAMP, CMMC, DORA, NIS2).
The thesaurus would require extension for new domains.
The benchmark size (58 requirements) limits statistical power; we compensate with 30 seeds and report confidence intervals throughout.
We plan to evaluate on real-world mapping data (e.g., the NIST SP 800-53 to ISO 27001 crosswalk published by NIST) in future work.


%% ================================================================
\section{Conclusion and Future Work}
\label{sec:conclusion}
%% ================================================================

We have presented the Regulatory Alignment Agent, a multi-step retrieval system for automated compliance mapping.
Through a controlled ablation study, we established three principal findings:
(1)~vocabulary mismatch between regulatory and implementation language is the primary retrieval bottleneck, and domain-aware query reformulation provides the largest gain ($\Delta\text{Top-1} = +0.116$, $p < 0.001$);
(2)~multi-backend fusion and cross-framework corroboration provide complementary improvements in ranking robustness and decision calibration;
(3)~the full agent significantly outperforms lexical baselines and a general-purpose neural encoder while producing auditable reasoning traces.

Several directions for future work are promising.
First, integrating domain-adapted language models (e.g., Legal-BERT~\cite{Chalkidis2020}, SecureBERT~\cite{Aghaei2022}) as additional retrieval backends within the agent's fusion framework.
Second, evaluating on real-world compliance mapping datasets, such as the NIST-published framework crosswalks.
Third, extending the agent with an LLM-based reasoning module that can generate natural-language explanations for mapping decisions, further improving auditability.
Fourth, investigating temporal aspects of compliance mapping, where regulations and controls evolve over time and previously valid mappings may become stale~\cite{Gama2014}.


%% ================================================================
\section*{Data Availability}
%% ================================================================
The benchmark (58 regulations, 110 controls, 86 mapping labels) is provided as Online Resource 1.

\section*{Code Availability}
The complete evaluation framework, all baselines, agent implementation, domain thesaurus, and ablation harness are provided as Online Resource 2 (Python CLI).
The semantic backend requires external model downloads (\texttt{pip install sentence-transformers}); all other components are self-contained.

\section*{Online Resources}
Online Resource 1: Hardened synthetic benchmark (CSV).\\
Online Resource 2: Reproducible evaluation program, agent implementation, and scripts to reproduce all tables.

\section*{Declarations}
\textbf{Funding:} This research received no external funding.\\
\textbf{Conflict of interest:} The author declares no competing interests.\\
\textbf{Ethics:} This study uses synthetic data and does not involve human subjects.

%% ================================================================
%% Bibliography
%% ================================================================

\begin{thebibliography}{99}

\bibitem{MarcusMaletic2003}
Marcus, A., Maletic, J.I.: Recovering documentation-to-source-code traceability links using latent semantic indexing. In: Proc.\ 25th International Conference on Software Engineering (ICSE), pp.\ 349--357 (2003). \textsc{doi}: 10.1109/ICSE.2003.1201194

\bibitem{Antoniol2002}
Antoniol, G., Canfora, G., Casazza, G., De Lucia, A., Merlo, E.: Recovering traceability links between code and documentation. IEEE Trans.\ Softw.\ Eng.\ \textbf{28}(10), 970--983 (2002). \textsc{doi}: 10.1109/TSE.2002.1041053

\bibitem{Borg2014}
Borg, M., Runeson, P., Ard\"o, A.: Recovering from a decade: a systematic mapping of information retrieval approaches to software traceability. Empir.\ Softw.\ Eng.\ \textbf{19}(6), 1565--1616 (2014). \textsc{doi}: 10.1007/s10664-013-9255-y

\bibitem{Cleland-Huang2014}
Cleland-Huang, J., Gotel, O.C.Z., Huffman Hayes, J., Mader, P., Zisman, A.: Software traceability: trends and future directions. In: Proc.\ Future of Software Engineering (FOSE), pp.\ 55--69 (2014). \textsc{doi}: 10.1145/2593882.2593891

\bibitem{McMillan2009}
McMillan, C., Poshyvanyk, D., Revelle, M.: Combining textual and structural analysis of software artifacts for traceability link recovery. In: Proc.\ TEFSE Workshop, pp.\ 41--48 (2009). \textsc{doi}: 10.1109/TEFSE.2009.5069582

\bibitem{Cuddeback2010}
Cuddeback, D., Dekhtyar, A., Huffman Hayes, J.: Automated requirements traceability: the study of human analysts. In: Proc.\ 18th IEEE RE, pp.\ 231--240 (2010). \textsc{doi}: 10.1109/RE.2010.35

\bibitem{Ali2013}
Ali, N., Gueheneuc, Y.-G., Antoniol, G.: Trustrace: Mining software repositories to improve the accuracy of requirement traceability links. IEEE Trans.\ Softw.\ Eng.\ \textbf{39}(5), 725--741 (2013). \textsc{doi}: 10.1109/TSE.2012.71

\bibitem{Oliveto2010}
Oliveto, R., Gethers, M., Poshyvanyk, D., De Lucia, A.: On the equivalence of information retrieval methods for automated traceability link recovery. In: Proc.\ 18th IEEE ICPC, pp.\ 68--71 (2010). \textsc{doi}: 10.1109/ICPC.2010.20

\bibitem{Guo2017}
Guo, J., Cheng, J., Cleland-Huang, J.: Semantically enhanced software traceability using deep learning techniques. In: Proc.\ 39th ICSE, pp.\ 3--14 (2017). \textsc{doi}: 10.1109/ICSE.2017.9

\bibitem{Devlin2019}
Devlin, J., Chang, M.-W., Lee, K., Toutanova, K.: BERT: Pre-training of deep bidirectional transformers for language understanding. In: Proc.\ NAACL-HLT, pp.\ 4171--4186 (2019). \textsc{doi}: 10.18653/v1/N19-1423

\bibitem{Chalkidis2020}
Chalkidis, I., Fergadiotis, M., Malakasiotis, P., Aletras, N., Androutsopoulos, I.: LEGAL-BERT: The Muppets straight out of Law School. In: Findings of EMNLP, pp.\ 2898--2904 (2020). \textsc{doi}: 10.18653/v1/2020.findings-emnlp.261

\bibitem{Aghaei2022}
Aghaei, E., Niu, X., Shadid, W., Al-Shaer, E.: SecureBERT: A domain-specific language model for cybersecurity. In: Proc.\ SecureComm (2022). arXiv:2204.02685

\bibitem{Robertson2009}
Robertson, S., Zaragoza, H.: The probabilistic relevance framework: BM25 and beyond. Found.\ Trends Inf.\ Retr.\ \textbf{3}(4), 333--389 (2009). \textsc{doi}: 10.1561/1500000019

\bibitem{Cormack2009}
Cormack, G.V., Clarke, C.L.A., Buettcher, S.: Reciprocal rank fusion outperforms condorcet and individual rank learning methods. In: Proc.\ 32nd ACM SIGIR, pp.\ 758--759 (2009). \textsc{doi}: 10.1145/1571941.1572114

\bibitem{Carpineto2012}
Carpineto, C., Romano, G.: A survey of automatic query expansion in information retrieval. ACM Comput.\ Surv.\ \textbf{44}(1), 1:1--1:50 (2012). \textsc{doi}: 10.1145/2071389.2071390

\bibitem{Rocchio1971}
Rocchio, J.J.: Relevance feedback in information retrieval. In: Salton, G.\ (ed.) The SMART Retrieval System: Experiments in Automatic Document Processing, pp.\ 313--323. Prentice Hall (1971)

\bibitem{Voorhees1994}
Voorhees, E.M.: Query expansion using lexical-semantic relations. In: Proc.\ 17th ACM SIGIR, pp.\ 61--69 (1994). \textsc{doi}: 10.1007/978-1-4471-2099-5\_7

\bibitem{Nogueira2019}
Nogueira, R., Cho, K.: Passage re-ranking with BERT. arXiv:1901.04085 (2019)

\bibitem{Aronson2010}
Aronson, A.R., Lang, F.-M.: An overview of MetaMap: historical perspective and recent advances. J.\ Am.\ Med.\ Inform.\ Assoc.\ \textbf{17}(3), 229--236 (2010). \textsc{doi}: 10.1136/jamia.2009.002733

\bibitem{Maxwell2017}
Maxwell, K.T., Schafer, B.: Concept and context in legal information retrieval. In: Proc.\ JURIX, pp.\ 63--72 (2017). \textsc{doi}: 10.3233/978-1-61499-838-9-63

\bibitem{Fox1994}
Fox, E.A., Shaw, J.A.: Combination of multiple searches. In: Proc.\ 2nd TREC, pp.\ 243--252 (1994)

\bibitem{ElYaniv2010}
El-Yaniv, R., Wiener, Y.: On the foundations of noise-free selective classification. J.\ Mach.\ Learn.\ Res.\ \textbf{11}, 1605--1641 (2010)

\bibitem{Geifman2017}
Geifman, Y., El-Yaniv, R.: Selective classification for deep neural networks. In: Proc.\ NeurIPS, pp.\ 4878--4887 (2017)

\bibitem{Geifman2019}
Geifman, Y., El-Yaniv, R.: SelectiveNet: A deep neural network with an integrated reject option. In: Proc.\ ICML, pp.\ 2151--2159 (2019)

\bibitem{Kamath2020}
Kamath, A., Jia, R., Liang, P.: Selective question answering under domain shift. In: Proc.\ ACL, pp.\ 5684--5696 (2020). \textsc{doi}: 10.18653/v1/2020.acl-main.503

\bibitem{Yao2023}
Yao, S., Zhao, J., Yu, D., Du, N., Shafran, I., Narasimhan, K., Cao, Y.: ReAct: Synergizing reasoning and acting in language models. In: Proc.\ ICLR (2023). arXiv:2210.03629

\bibitem{Yang2024}
Yang, J., Jimenez, C.E., Wettig, A., Liber, K., Narasimhan, K., Press, O.: SWE-agent: Agent-computer interfaces enable automated software engineering. arXiv:2405.15793 (2024)

\bibitem{Bran2024}
Bran, A.M., Cox, S., Schilter, O., Baldassari, C., White, A.D., Schwaller, P.: Augmenting large language models with chemistry tools. Nat.\ Mach.\ Intell.\ \textbf{6}, 525--535 (2024). \textsc{doi}: 10.1038/s42256-024-00832-8

\bibitem{Shao2024}
Shao, Z., Gong, Y., Shen, Y., Huang, M., Duan, N., Chen, W.: Assisting in writing Wikipedia-like articles from scratch with large language models. In: Proc.\ NAACL, pp.\ 6252--6278 (2024). \textsc{doi}: 10.18653/v1/2024.naacl-long.347

\bibitem{Reimers2019}
Reimers, N., Gurevych, I.: Sentence-BERT: Sentence embeddings using Siamese BERT-networks. In: Proc.\ EMNLP, pp.\ 3982--3992 (2019). \textsc{doi}: 10.18653/v1/D19-1410

\bibitem{Lample2018}
Lample, G., Conneau, A., Denoyer, L., Ranzato, M.: Unsupervised machine translation using monolingual corpora only. In: Proc.\ ICLR (2018). arXiv:1711.00043

\bibitem{Breaux2008}
Breaux, T.D., Ant\'on, A.I.: Analyzing regulatory rules for privacy and security requirements. IEEE Trans.\ Softw.\ Eng.\ \textbf{34}(1), 5--20 (2008). \textsc{doi}: 10.1109/TSE.2007.70746

\bibitem{Gama2014}
Gama, J., \v{Z}liobait\.{e}, I., Bifet, A., Pechenizkiy, M., Bouchachia, A.: A survey on concept drift adaptation. ACM Comput.\ Surv.\ \textbf{46}(4), 44:1--44:37 (2014). \textsc{doi}: 10.1145/2523813

\end{thebibliography}

\end{document}
